% numerical_validation.tex —— scuc 数值验证证据
% 规范：docs/modules/README.md 数值证据规范——只引用真实存在的测试目标、
% 基线文件与可复现命令；未执行的交叉验证不得写成已完成。

\section{数值验证}

\subsection{测试目标注册（CMake 事实）}

根 \texttt{CMakeLists.txt} 注册两个覆盖本模块的测试可执行文件（均要求
\texttt{MIPSOLVERS\_BUILD\_SCUC=ON}，默认开；运行时目录为 \texttt{tests/}）：

\begin{center}
\begin{footnotesize}
\begin{tabular}{@{}llll@{}}
\toprule
\textbf{目标} & \textbf{源文件} & \textbf{ctest 标签} & \textbf{说明}\\
\midrule
\texttt{test\_scuc\_module} & \srcpath{tests/test\_scuc\_module.cpp} &
\texttt{integration} & 模块级单元/集成测试\\
\texttt{test\_market\_simulation} & \srcpath{tests/test\_market\_simulation.cpp} &
\texttt{benchmark}（\texttt{TIMEOUT 600}） & 多求解器基准与物理一致性\\
\bottomrule
\end{tabular}
\end{footnotesize}
\end{center}

两者都直接编译 \srcpath{src/scuc/scuc.cpp}（后者还编译
\srcpath{src/scuc/case\_builder.cpp}）并链接 \texttt{mipsolvers} 与 Catch2。
复现命令（在已配置的构建目录下）：

\begin{verbatim}
ctest -R test_scuc_module --output-on-failure
ctest -R test_market_simulation --output-on-failure
# 或直接运行（Catch2 标签过滤）：
./tests/test_scuc_module "[scuc]"
./tests/test_market_simulation "[market][physics]"
\end{verbatim}

\subsection{覆盖矩阵：\texttt{test\_scuc\_module}}

\srcpath{tests/test\_scuc\_module.cpp} 以 2 母线、2 机组、3 时段、限额 150 MW
单支路的小问题（\srcpath{tests/test\_scuc\_module.cpp:make\_small\_problem\_json}，
\fld{mip\_gap}=0.005、时限 60 s）覆盖：

\begin{center}
\begin{footnotesize}
\begin{tabular}{@{}ll@{}}
\toprule
\textbf{TEST\_CASE（标签）} & \textbf{断言内容}\\
\midrule
JSON round-trip（\texttt{[scuc][unit]}） &
解析无异常；机组/支路/负荷数、\fld{num\_periods}、\fld{num\_buses}、
分段报价读回正确\\
small 2-bus problem converges（\texttt{[scuc][integration]}） &
SCUC 收敛、目标 $>0$、组合矩阵形状 $[2][3]$、切负荷 $|P^{shed}|<1$ MW、
时段 0 功率平衡（容差 5 MW）、能量/总成本 $>0$\\
SCED LP refines dispatch（\texttt{[scuc][integration]}） &
SCED 收敛且组合与 SCUC 一致（固定二进制口径，容差 0.01）\\
LMP values are physically consistent（\texttt{[scuc][integration]}） &
$0\le$ avg LMP $\le 300$ \$/MWh、min $\le$ avg $\le$ max、矩阵形状
$[n_b][T]$、高负荷时段 LMP 不低于低负荷时段（容差 0.1）\\
output JSON serialisation round-trips（\texttt{[scuc][unit]}） &
输出 JSON 可解析、含 \texttt{scuc}/\texttt{meta}/\texttt{lmp} 块且
收敛标志一致\\
\bottomrule
\end{tabular}
\end{footnotesize}
\end{center}

\subsection{覆盖矩阵：\texttt{test\_market\_simulation}}

\srcpath{tests/test\_market\_simulation.cpp} 基于 case\_builder 算例组织
（基准族对所有\textbf{当前构建可用}的 MILP 后端逐一运行，
\srcpath{tests/test\_market\_simulation.cpp:available\_milp\_solvers}；单后端
失败不致命、只记录，Auto 模式必须收敛）：

\begin{center}
\begin{footnotesize}
\begin{tabular}{@{}ll@{}}
\toprule
\textbf{TEST\_CASE（标签）} & \textbf{内容}\\
\midrule
3-bus 全求解器性能基准 T=6（\texttt{[market][benchmark]}） &
各后端测速对比；Auto 必须收敛\\
6-bus 风电+储能 全求解器基准 T=8（\texttt{[market][benchmark]}） &
含 \fld{use\_binary\_indicators} 之外全部储能路径（$M_2=50$）\\
IEEE 39-bus 快速基准 T=4 / 24h 风光 / 24h 全资源
（\texttt{[market][benchmark][ieee39]}） & 39 母线规模收敛与用时\\
备用约束验证（\texttt{[market][physics]}） &
$r^{spin}=0.10$、$r^{up}=r^{dn}=0.05$、$r^{neg}=0.10$、$R^{pfr}=10$ MW
下逐时段旋转备用 $\ge r^{spin}D_t-1$ MW\\
Gantt/出力/LMP 热图/SOC 时序（\texttt{[market][viz]}） &
ASCII 可视化 + 物理一致性抽查\\
IEEE 39-bus 24h 完整流程（\texttt{[market][viz][ieee39]}） &
三阶段全跑通 + JSON 导出\\
MIP Gap 收敛速度分析（\texttt{[market][benchmark]}） & 间隙-时间曲线记录\\
StrictHiGHS 39-bus root IPM + crossover 正确性
（\texttt{[market][strict\_highs][ipm\_root][ieee39]}） &
经 \srcpath{src/scuc/scuc.cpp:build\_scuc\_mip} 取裸模型直调
\texttt{solve\_milp\_bc}：IPM 根 + crossover 与单纯形根目标相对差 $<0.5\%$\\
StrictHiGHS 纯 IPM（无 crossover）（\texttt{...},\texttt{[pure\_ipm]}） &
记录无暖启动的节点 LP 代价；与 IPM+crossover 目标差 $<0.5\%$\\
\bottomrule
\end{tabular}
\end{footnotesize}
\end{center}

\subsection{CLI 冒烟用法}

构建后（可执行文件在构建输出目录；安装后经 \texttt{bin/}）：

\begin{verbatim}
# 生成算例 JSON
scuc_case_builder --case 3bus --T 6 --output case_3bus.json
scuc_case_builder --case ieee39 --T 24 --wind --solar --output case_39.json
# 求解（stderr 打印配置摘要与各阶段收敛状态）
scuc_solve case_3bus.json result_3bus.json --solver Auto --indent 2
scuc_solve case_39.json  result_39.json  --solver StrictHiGHS
# 仅 SCUC 阶段：--no-sced --no-lmp
\end{verbatim}

\texttt{scuc\_solve} 在 SCUC 未收敛时返回非零退出码（
\srcpath{src/scuc/main.cpp:main}），可直接接入调度脚本做失败告警；
\texttt{--solver} 覆盖 JSON 中的 \fld{solver} 字段。

\subsection{已知覆盖边界}

\begin{boundarynote}
\begin{itemize}
  \item 上述测试覆盖：JSON 往返、小规模收敛与功率平衡、SCED 组合一致性、
        LMP 数量级与单调性、备用需求、多后端基准、StrictHiGHS 根 LP 内核
        交叉校核（0.5\% 容差）；
  \item \textbf{未覆盖}：IEEE 118 母线规模无 \texttt{tests/} 注册测试
        （其求解器对比由 \srcpath{benchmark/bench\_01\_ieee118\_solvers.py}
        等 benchmark 脚本承载，不属于 ctest 套件）；无与外部参考求解器
        （如商业 MILP）在本仓库内的系统性目标值比对记录；N-1、多场景、
        交流潮流无实现亦无测试；断面（\fld{sections}）、机组群
        （\fld{generator\_groups}）、直流线、二进制储能模式
        （\fld{use\_binary\_indicators}）、热/温/冷启动字段无专门断言；
  \item 基准族对单后端失败\textbf{宽容}（仅记录），只有 Auto 模式的收敛是
        硬断言；阅读基准输出时须区分"求解器不可用"与"求解失败"；
  \item 容差口径：功率平衡 5 MW、LMP 单调性 0.1 \$/MWh、IPM/单纯形目标
        比对 0.5\%，均为小算例口径，不外推到大算例。
\end{itemize}
\end{boundarynote}
